\documentclass[11pt]{article}

\usepackage[margin=1.05in]{geometry}
\usepackage[T1]{fontenc}
\usepackage{lmodern}
\usepackage{amsmath,amssymb,amsthm,mathtools}
\usepackage{enumitem}
\usepackage{booktabs}
\usepackage{microtype}
\usepackage[colorlinks=true,linkcolor=blue,citecolor=blue,urlcolor=blue]{hyperref}

\newtheorem{theorem}{Theorem}[section]
\newtheorem{proposition}[theorem]{Proposition}
\newtheorem{lemma}[theorem]{Lemma}
\newtheorem{corollary}[theorem]{Corollary}
\newtheorem{inputtheorem}[theorem]{Input theorem}
\newtheorem{criterion}[theorem]{Criterion}
\theoremstyle{definition}
\newtheorem{definition}[theorem]{Definition}
\theoremstyle{remark}
\newtheorem{remark}[theorem]{Remark}
\newtheorem{problem}[theorem]{Problem}

\newcommand{\Pp}{\mathbb P}
\newcommand{\Ee}{\mathbb E}
\newcommand{\Var}{\operatorname{Var}}
\newcommand{\Cov}{\operatorname{Cov}}
\newcommand{\one}{\mathbf 1}
\newcommand{\dd}{\,d}
\newcommand{\Exp}{\operatorname{Exp}}
\newcommand{\Pois}{\operatorname{Pois}}
\newcommand{\ord}{\operatorname{ord}}
\newcommand{\lcm}{\operatorname{lcm}}
\newcommand{\Cyc}{\operatorname{Cyc}}
\newcommand{\TV}{\mathrm{TV}}
\newcommand{\cL}{\mathcal L}
\newcommand{\cO}{\mathcal O}
\newcommand{\cI}{\mathcal I}
\newcommand{\cD}{\mathcal D}
\newcommand{\cA}{\mathcal A}
\newcommand{\cG}{\mathcal G}
\newcommand{\R}{\mathbb R}
\newcommand{\Z}{\mathbb Z}

\title{Functional Erd\H{o}s--Tur\'an Universality for\\
Comparable-Rate Luce Permutations}
\author{Christopher D. Long}
\date{Working branch --- functional and multivariate strengthening}

\begin{document}
\maketitle

\begin{abstract}
Let $\sigma_n$ be the Luce permutation induced by independent exponential
clocks whose rates, after a common rescaling, lie in $[1,\kappa]$ for one
fixed $\kappa<\infty$.  This working branch strengthens the scalar
Erd\H{o}s--Tur\'an theorem to a process-level universality statement.  If
$C_{n,r}$ is the number of $r$-cycles, set
\[
 K_n(t)=\sum_{r\le n^t}C_{n,r},
 \qquad
 \cO_n(t)=\log\lcm\{r\le n^t:C_{n,r}>0\},
 \qquad 0\le t\le1.
\]
We prove, conditional only on the switching inputs already isolated in the
scalar paper, that
\[
 \left(
 \frac{K_n(t)-t\log n}{\sqrt{\log n}},
 \frac{\cO_n(t)-\frac12t^2(\log n)^2}
      {\sqrt{\frac13(\log n)^3}}
 \right)_{0\le t\le1}
 \Rightarrow
 \left(
 W(t),\sqrt3\int_0^t u\,\dd W(u)
 \right)_{0\le t\le1}
\]
in $D([0,1],\R^2)$.  The main new intermediate statement is a
high-dimensional product-Poisson approximation: throughout the growing
window
\[
 (\log n)^3<r\le \frac{n}{(\log n)^{A_0}},
\]
the entire cycle-count vector is asymptotically independent with coordinates
$\Pois(1/r)$ in total variation.  It follows from an approximate Campbell
identity obtained directly from the marked one-switch formula and a
configuration-space Poisson Stein argument.  The existing three-clock and
von Mangoldt estimates then transfer the additive process to the truncated
least-common-multiple process uniformly in time.  We also derive an
isonormal Gaussian limit for bounded logarithmic-scale cycle weights and the
joint endpoint law
\[
 \left(
 \frac{K_n-\log n}{\sqrt{\log n}},
 \frac{\log\ord(\sigma_n)-\frac12(\log n)^2}
      {\sqrt{\frac13(\log n)^3}}
 \right)
 \Rightarrow
 N\!\left(0,
 \begin{pmatrix}1&\sqrt3/2\\ \sqrt3/2&1\end{pmatrix}\right).
\]
The final section records how this qualitative process proof interfaces with,
but does not replace, the frequency-sensitive Berry--Esseen branch.
\end{abstract}

\tableofcontents

\section{Introduction and theorem map}

For a permutation $\pi\in S_n$,
\[
 \ord(\pi)=\lcm\{\text{cycle lengths of }\pi\}.
\]
For a uniform permutation, Erd\H{o}s and Tur\'an proved that
$\log\ord(\pi)$ is asymptotically Gaussian after centering by
$\frac12(\log n)^2$ and scaling by $\sqrt{\frac13(\log n)^3}$
\cite{ErdosTuran1967}.  The scalar comparable-rate Luce theorem proves that
this normalization survives arbitrary bounded triangular arrays of clock
rates; see \cite{LongOrder2026}.  The purpose of the present branch is to
identify the full logarithmic-scale Gaussian field behind that endpoint
law.

Throughout, $W$ denotes a standard Brownian motion.  Define
\begin{equation}\label{eq:processes-intro}
 K_n(t)=\sum_{r\le n^t}C_{n,r},
 \qquad
 B_n(t)=\sum_{r\le n^t}(\log r)C_{n,r},
\end{equation}
and
\begin{equation}\label{eq:order-process-intro}
 \cO_n(t)=\log\lcm\{r\le n^t:C_{n,r}>0\}.
\end{equation}
The convention is that the least common multiple of the empty set is one.

\begin{theorem}[Bivariate functional Erd\H{o}s--Tur\'an theorem]
\label{thm:main-functional}
Let $\sigma_n$ be a Luce permutation whose positive rates satisfy
\[
 \sup_n\frac{\max_i\theta_{n,i}}{\min_i\theta_{n,i}}<\infty.
\]
Then, in $D([0,1],\R^2)$ with the $J_1$ topology,
\begin{equation}\label{eq:main-functional}
 \boxed{
 \left(
 \frac{K_n(t)-t\log n}{\sqrt{\log n}},
 \frac{\cO_n(t)-\frac12t^2(\log n)^2}
      {\sqrt{\frac13(\log n)^3}}
 \right)_{0\le t\le1}
 \Rightarrow
 \left(
 W(t),\sqrt3\int_0^t u\,\dd W(u)
 \right)_{0\le t\le1}.}
\end{equation}
The convergence is uniform over all triangular arrays satisfying a fixed
comparability bound.
\end{theorem}

The limiting covariance kernels are
\begin{align}
 \Cov(W(s),W(t))&=s\wedge t,\label{eq:cov-count-count}\\
 \Cov\left(W(s),\sqrt3\int_0^t u\,\dd W(u)\right)
 &=\frac{\sqrt3}{2}(s\wedge t)^2,\label{eq:cov-count-order}\\
 \Cov\left(\sqrt3\int_0^s u\,\dd W(u),
            \sqrt3\int_0^t u\,\dd W(u)\right)
 &=(s\wedge t)^3.\label{eq:cov-order-order}
\end{align}
In particular, the second coordinate alone has the same law as
$B(t^3)$ for another standard Brownian motion $B$.

\begin{corollary}[Joint endpoint law]
\label{cor:endpoint-joint}
Writing $K_n=K_n(1)$,
\begin{equation}\label{eq:endpoint-joint}
 \boxed{
 \left(
 \frac{K_n-\log n}{\sqrt{\log n}},
 \frac{\log\ord(\sigma_n)-\frac12(\log n)^2}
      {\sqrt{\frac13(\log n)^3}}
 \right)
 \Rightarrow
 N\!\left(0,
 \begin{pmatrix}
 1&\sqrt3/2\\[1mm]
 \sqrt3/2&1
 \end{pmatrix}\right).}
\end{equation}
\end{corollary}

The process theorem follows from a stronger intermediate-window statement.

\begin{theorem}[Product-Poisson approximation of the intermediate cycle field]
\label{thm:TV-main}
After the normalization $1\le\theta_{n,i}\le\kappa$, choose the cutoffs
$b_n,U_n$ in \eqref{eq:cutoffs}.  Let
\[
 \xi_n=(C_{n,r})_{b_n<r\le U_n}.
\]
Let $(Z_{n,r})_{b_n<r\le U_n}$ be independent with
\[
 Z_{n,r}\sim\Pois(1/r).
\]
Then
\begin{equation}\label{eq:TV-main}
 \boxed{
 d_{\TV}\left(
 \cL(\xi_n),
 \bigotimes_{b_n<r\le U_n}\Pois(1/r)
 \right)\longrightarrow0.}
\end{equation}
The convergence is uniform over all rate arrays in $[1,\kappa]$.
\end{theorem}

This theorem is stronger than what is required for a scalar characteristic
function.  It says that the entire growing cycle vector in the range that
carries the Erd\H{o}s--Tur\'an variance is asymptotically the classical
independent Poisson field.

There is also a convenient Gaussian-random-measure formulation.  Put
$\ell_n=\log n$ and $u_{n,r}=\log r/\ell_n$.

\begin{theorem}[Logarithmic-scale Gaussian cycle field]
\label{thm:isonormal}
For every fixed collection of bounded Riemann-integrable functions
$f_1,\ldots,f_d:[0,1]\to\R$, define
\begin{equation}\label{eq:Gn-f}
 \mathbb G_n(f)
 =\frac1{\sqrt{\ell_n}}
 \left[
 \sum_{r=1}^n f(u_{n,r})C_{n,r}
 -\sum_{r=1}^n\frac{f(u_{n,r})}{r}
 \right].
\end{equation}
Then
\begin{equation}\label{eq:isonormal}
 \bigl(\mathbb G_n(f_1),\ldots,\mathbb G_n(f_d)\bigr)
 \Rightarrow N(0,\Sigma),
 \qquad
 \Sigma_{ij}=\int_0^1f_i(u)f_j(u)\,\dd u.
\end{equation}
Equivalently,
\[
 \mathbb G_n(f)\Rightarrow\int_0^1 f(u)\,\dd W(u)
\]
in every fixed finite family of test functions.
\end{theorem}

The indicator weights $f(u)=\one_{[0,t]}(u)$ give the cycle-count process.
The weights $f(u)=\sqrt3u\one_{[0,t]}(u)$ give the normalized additive order
process.  The least-common-multiple process then follows from a uniform
additive-to-LCM transfer.

\subsection*{Proof architecture}

The argument has five modules.

\begin{enumerate}[label=\textnormal{(\roman*)},leftmargin=2.4em]
\item The quantitative endpoint theory from the scalar paper gives a
one-block bound for
\[
 \Delta_{n,r}=\Ee\left|S_{n,r}/n-1\right|.
\]
Repeating its dyadic summation without the logarithmic weight yields
\[
 \sum_j\sup_{2^jb_n\le r<2^{j+1}b_n}\Delta_{n,r}=o(1).
\]

\item The marked one-switch identity is tested against arbitrary bounded
functions of the complete intermediate cycle vector.  Away from the short
source-cycle exception, the switch adds exactly one $r$-cycle and changes no
other intermediate coordinate.  This gives an approximate Campbell identity.

\item An immigration--death Stein equation converts the Campbell identity
into total-variation approximation by a product of Poisson variables.  The
Stein constant is dimension-free.

\item The product-Poisson field satisfies a standard triangular-array
martingale functional CLT on the logarithmic coordinate
$u=\log r/\log n$.  Total variation transfers this process theorem to the
Luce cycle vector.  Crude cycle bounds restore the omitted short and long
ranges.

\item The discrepancy $B_n(t)-\cO_n(t)$ is nonnegative and nondecreasing in
$t$.  Hence its path supremum equals its endpoint value.  The existing
three-clock/von Mangoldt estimate
\[
 \Ee[B_n(1)-\cO_n(1)]=O_\kappa(\log n\log\log n)
\]
therefore transfers the whole additive process to the LCM process.
\end{enumerate}

\subsection*{Status of this branch}

The draft is written as a proof conditional on the same theorem-level inputs
already used by the scalar paper.  The only genuinely new probabilistic step
is Theorem~\ref{thm:TV-main}.  The branch deliberately keeps the scalar
characteristic-function proof and the Berry--Esseen supplement separate until
the product-Poisson argument has been independently audited.

\section{Model and imported switching inputs}

After a common rescaling, assume throughout that
\begin{equation}\label{eq:comparable}
 1\le\theta_{n,i}\le\kappa,
 \qquad 1\le i\le n,
\end{equation}
for one fixed $\kappa<\infty$.  Let
\[
 X_{n,i}\sim\Exp(\theta_{n,i}),
 \qquad 1\le i\le n,
\]
be independent, and let $\sigma_n(i)$ be the rank of $X_{n,i}$ among the
$n$ clocks.  Write $C_{n,r}$ for the number of $r$-cycles and $L_n(i)$ for
the cycle length containing $i$.

For distinct $i,j$, put
\begin{equation}\label{eq:Rij}
 R_{ij}=\exp\{-(\theta_{n,i}-\theta_{n,j})(X_{n,j}-X_{n,i})\},
 \qquad R_{ii}=1,
\end{equation}
and define
\begin{equation}\label{eq:S-n-r}
 S_{n,r}=\sum_{i:L_n(i)>r}R_{i,\sigma_n^r(i)}.
\end{equation}

\begin{inputtheorem}[Marked one-switch identity, moments, and tails]
\label{input:one-switch}
For every $1\le r<n$ and every bounded complex function
$\Phi:S_n\to\mathbb C$,
\begin{equation}\label{eq:marked-switch}
 r(n-r)\Ee[C_{n,r}\Phi(\sigma_n)]
 =\Ee\sum_{i:L_n(i)>r}
 R_{i,\sigma_n^r(i)}
 \Phi\bigl(\sigma_n\circ(i\,\sigma_n^r(i))\bigr).
\end{equation}
In particular,
\begin{equation}\label{eq:mean-switch}
 r(n-r)\Ee C_{n,r}=\Ee S_{n,r}.
\end{equation}
Put
\[
 \lambda_\kappa=
 \begin{cases}
 \infty,&\kappa=1,\\[1mm]
 \kappa/(\kappa-1),&\kappa>1.
 \end{cases}
\]
For every $1<p<\lambda_\kappa$,
\begin{equation}\label{eq:selected-p}
 \sup_{n,r}\frac1n\Ee\sum_{i:L_n(i)>r}
 R_{i,\sigma_n^r(i)}^p<\infty.
\end{equation}
Moreover, the truncation and fresh-pair tail estimates stated in the scalar
paper hold uniformly, and $\Ee R_{ij}=1$ for every $i,j$.
\end{inputtheorem}

This is the exact input used in the current scalar draft
\cite{LongOrder2026}.  It implies the following crude bounds.

\begin{lemma}[Crude cycle, root, and tail bounds]
\label{lem:crude}
Uniformly in $1\le r<n$,
\begin{equation}\label{eq:Cr-crude}
 \Ee C_{n,r}\le \frac{C_\kappa n}{r(n-r)}.
\end{equation}
If $I_n$ is uniform on $[n]$, independently of the clocks, then for
$a\le n/2$,
\begin{equation}\label{eq:root-crude}
 \Pp\{L_n(I_n)\le a\}\le C_\kappa a/n.
\end{equation}
Consequently, for $b\le n/2$ and $U\le n/2$,
\begin{align}
 \Ee\sum_{r\le b}C_{n,r}&=O_\kappa(\log(eb)),
 \label{eq:small-count-tail}\\
 \Ee\sum_{r>U}C_{n,r}&=O_\kappa(1+\log(n/U)),
 \label{eq:large-count-tail}\\
 \Ee\sum_{r\le b}(\log r)C_{n,r}
 &=O_\kappa((\log(eb))^2),
 \label{eq:small-weight-tail}\\
 \Ee\sum_{r>U}(\log r)C_{n,r}
 &\le C_\kappa\log n\{1+\log(n/U)\}.
 \label{eq:large-weight-tail}
\end{align}
\end{lemma}

\begin{proof}
Equation~\eqref{eq:Cr-crude} follows from
\eqref{eq:mean-switch} and $\Ee S_{n,r}\le C_\kappa n$.  Since
\[
 \Pp\{L_n(I_n)=r\}=\frac{r\Ee C_{n,r}}n,
\]
summing over $r\le a$ proves \eqref{eq:root-crude}.  For $r\le n/2$,
\eqref{eq:Cr-crude} gives $\Ee C_{n,r}\le C_\kappa/r$.  Harmonic summation
proves the first and third estimates and the portions of the second and
fourth estimates below $n/2$.  There is deterministically at most one cycle
longer than $n/2$.
\end{proof}

Fix once and for all
\begin{equation}\label{eq:p-beta}
 1<p<\lambda_\kappa,
 \qquad
 \beta=1-\frac1p>0.
\end{equation}
The process argument needs a slightly stronger upper cutoff than the scalar
CLT.  To align this branch with the Berry--Esseen supplement, choose
\begin{equation}\label{eq:A0-choice}
 A_0>\max\left\{32\kappa^2,\frac4\beta,4\right\},
\end{equation}
and put
\begin{equation}\label{eq:cutoffs}
 b_n=\left\lceil(\log n)^3\right\rceil,
 \qquad
 U_n=\left\lfloor\frac{n}{(\log n)^{A_0}}\right\rfloor.
\end{equation}
For all sufficiently large $n$,
\begin{equation}\label{eq:cutoff-range}
 2U_n\le n/2.
\end{equation}

Put
\begin{equation}\label{eq:Delta}
 \Delta_{n,r}=\Ee\left|\frac{S_{n,r}}n-1\right|.
\end{equation}
For a dyadic lower endpoint $a$, set
\begin{equation}\label{eq:block-parameters}
 \delta=\frac{4a}{n},
 \qquad
 L_a=\log\frac1\delta,
 \qquad
 m=\left\lfloor\min\left\{\frac a3,\sqrt{L_a}\right\}\right\rfloor.
\end{equation}

\begin{inputtheorem}[Quantitative one-block switching]
\label{input:block}
There is $C_\kappa<\infty$ such that, uniformly for
$b_n\le a\le U_n$ and $a\le r<2a$,
\begin{equation}\label{eq:block-delta}
 \Delta_{n,r}
 \le C_\kappa(\Gamma_{n,a}+n^{-1})^{\eta},
 \qquad
 \eta=\frac1{2\kappa},
\end{equation}
where
\begin{align}
 \Gamma_{n,a}\le C\bigg[&
 (1-\alpha_0)^{m-1}
 +m\delta^{1/4}
 +\frac{m^2}{n\delta^{1/2}}
 +m\delta^{1/(4\kappa)}
 +\frac{m\log n}{n\delta^{1/2}}
 \notag\\
 &+\delta^{1/(4\kappa)}
 +\delta^{1/2}L_a
 +\delta
 +g_{n,a}\bigg],
 \label{eq:Gamma}
\end{align}
with
\begin{equation}\label{eq:g-na}
 g_{n,a}\le C\left[
 \delta^{-1/4}L_a e^{-cn\delta^{1/2}}
 +n^2\delta^{-1/4}e^{-cn\delta^{3/4}}
 +n^{-2}
 \right].
\end{equation}
All constants depend only on $\kappa$.
\end{inputtheorem}

This theorem is the quantitative finite-endpoint interface proved and used in
the current scalar paper.  We also import the arithmetic output of its
three-clock section.

\begin{inputtheorem}[Two-cycle domination and endpoint product--LCM bound]
\label{input:arithmetic}
Whenever $r+s\le n/2$,
\begin{equation}\label{eq:pair-input}
 \Ee\left[C_{n,r}(C_{n,s}-\one_{\{r=s\}})\right]
 \le\frac{C_\kappa}{rs}.
\end{equation}
Moreover, with
\[
 B_n=B_n(1)=\sum_{r=1}^n(\log r)C_{n,r},
 \qquad O_n=\ord(\sigma_n),
\]
one has
\begin{equation}\label{eq:product-lcm-input}
 \Ee[B_n-\log O_n]
 \le C_\kappa\log n\log\log n.
\end{equation}
\end{inputtheorem}

The pair estimate follows from the exact rate-oriented three-clock identity;
the endpoint discrepancy follows by the von Mangoldt overlap decomposition.
Both are proved in \cite{LongOrder2026}.

\section{Unweighted dyadic summability}

The scalar proof sums $\Delta_{n,r}$ with a logarithmic block weight.  The
Campbell identity needs the stronger unweighted sum to tend to zero.
Let
\begin{equation}\label{eq:dyadic-aj}
 a_j=2^jb_n,
\end{equation}
and sum over those $j$ for which $a_j\le U_n$.

\begin{proposition}[Unweighted dyadic switching summability]
\label{prop:unweighted}
One has
\begin{equation}\label{eq:unweighted}
 \boxed{
 \sum_{j:a_j\le U_n}
 \sup_{a_j\le r<2a_j}\Delta_{n,r}
 \longrightarrow0.}
\end{equation}
The convergence is uniform over all rate arrays in $[1,\kappa]$.
\end{proposition}

\begin{proof}
Write
\[
 \delta_j=\frac{4a_j}{n},
 \qquad
 L_j=\log\frac1{\delta_j}.
\]
Since $a_j\ge b_n=(\log n)^3$ and $L_j\le\log n$, the minimum in
\eqref{eq:block-parameters} is attained by $\sqrt{L_j}$ for all large $n$,
uniformly in $j$.  Thus $m_j\asymp\sqrt{L_j}$.  Also
\begin{equation}\label{eq:Lj-min}
 L_j\ge A_0\log\log n-O(1).
\end{equation}
Because $0<\eta\le1$, the map $x\mapsto x^\eta$ is subadditive.  We may
therefore sum the contributions of the terms in \eqref{eq:Gamma}
separately.

For the reset term,
\begin{align*}
 \sum_j(1-\alpha_0)^{\eta(m_j-1)}
 &\le C\sum_{L_j\ge A_0\log\log n-O(1)}e^{-c\sqrt{L_j}}\\
 &\le C(1+\sqrt{\log\log n})
 e^{-c'\sqrt{\log\log n}}
 \longrightarrow0.
\end{align*}
Here the $L_j$ lie on an arithmetic lattice with fixed spacing.

After taking the $\eta$-power, every term involving a positive power of
$\delta_j$ is bounded by
\[
 C L_j^d e^{-\gamma L_j}
\]
for fixed $d\ge0$ and $\gamma>0$.  The smallest exponent comes from
$m\delta^{1/(4\kappa)}$ and is at least $1/(8\kappa^2)$.  Hence
\begin{equation}\label{eq:delta-sum}
 \sum_j L_j^d e^{-\gamma L_j}
 \le C(\log\log n)^d(\log n)^{-\gamma A_0}
 \longrightarrow0.
\end{equation}

For the terms containing $(n\delta_j^{1/2})^{-1}$, use
$m_j\le\sqrt{\log n}$ and $a_j\ge(\log n)^3$ to obtain
\[
 \frac{m_j^2+m_j\log n}{n\delta_j^{1/2}}
 \le C\frac{(\log n)^{3/2}}{\sqrt{na_j}}
 \le \frac{C}{\sqrt n}.
\]
After taking the $\eta$-power and summing over $O(\log n)$ blocks, the
contribution is $O((\log n)n^{-\eta/2})=o(1)$.

The first two terms in \eqref{eq:g-na} are superpolynomially small uniformly
for $a_j\ge(\log n)^3$.  The $n^{-2}$ term and the additional $n^{-1}$ in
\eqref{eq:block-delta} contribute $O((\log n)n^{-\eta})=o(1)$.  Combining
the estimates proves \eqref{eq:unweighted}.
\end{proof}

\begin{remark}[A usable explicit error]
\label{rem:explicit-unweighted}
The proof gives a deterministic $\varepsilon_n^{\mathrm{sw}}\downarrow0$,
depending only on $\kappa$, such that
\[
 \sum_j\sup_{a_j\le r<2a_j}\Delta_{n,r}
 \le\varepsilon_n^{\mathrm{sw}}.
\]
One may take a sum of terms of the schematic form
\[
 (1+\sqrt{\log\log n})e^{-c_\kappa\sqrt{\log\log n}},
 \quad
 (\log\log n)^d(\log n)^{-A_0/(8\kappa^2)},
 \quad
 (\log n)n^{-1/(4\kappa)}.
\]
No optimization is needed for the qualitative process theorem.
\end{remark}

\section{An approximate Campbell identity}

Let
\begin{equation}\label{eq:I-n}
 \cI_n=\{r\in\Z:b_n<r\le U_n\},
 \qquad
 \xi_n=(C_{n,r})_{r\in\cI_n}\in\Z_+^{\cI_n}.
\end{equation}
For $r\in\cI_n$, let $e_r$ denote the corresponding unit vector and put
\begin{equation}\label{eq:a-n-r}
 a_{n,r}=\frac{n}{r(n-r)}=\frac1r+\frac1{n-r}.
\end{equation}

The switch in \eqref{eq:marked-switch} cuts a single source cycle at two
marked points.

\begin{lemma}[Exact intermediate-vector effect of a good split]
\label{lem:good-split}
Fix $r\in\cI_n$ and a source label $i$ with
\[
 L_n(i)>r+U_n.
\]
Then
\begin{equation}\label{eq:good-split-vector}
 \xi_n\bigl(\sigma_n\circ(i\,\sigma_n^r(i))\bigr)
 =\xi_n(\sigma_n)+e_r.
\end{equation}
\end{lemma}

\begin{proof}
The transposition cuts the source cycle of length $L=L_n(i)$ into cycles of
lengths $r$ and $L-r$.  The source length satisfies $L>U_n$, and the residual
length satisfies $L-r>U_n$.  Thus the source and residual lengths are both
outside $\cI_n$, while the new $r$-cycle lies in $\cI_n$.  All other cycles
are unchanged.
\end{proof}

The exceptional source-cycle probability is uniformly summable.

\begin{lemma}[Exceptional source-cycle estimate]
\label{lem:exceptional}
Uniformly for $r\in\cI_n$,
\begin{equation}\label{eq:exceptional}
 \frac1n\Ee\sum_{i:L_n(i)>r}
 R_{i,\sigma_n^r(i)}
 \one_{\{L_n(i)\le r+U_n\}}
 \le C_\kappa\left(\frac{U_n}{n}\right)^\beta.
\end{equation}
\end{lemma}

\begin{proof}
Let $I_n$ be uniform on $[n]$, independently of the clocks.  The left side
is
\[
 \Ee\left[
 R_{I_n,\sigma_n^r(I_n)}
 \one_{\{r<L_n(I_n)\le r+U_n\}}
 \right].
\]
Apply H\"older's inequality with exponents $p$ and $p/(p-1)$.  The selected
$p$th moment is uniformly bounded by \eqref{eq:selected-p}.  Since
$r+U_n\le2U_n\le n/2$, \eqref{eq:root-crude} gives
\[
 \Pp\{L_n(I_n)\le r+U_n\}
 \le C_\kappa U_n/n.
\]
The conjugate H\"older exponent contributes the power
$\beta=1-1/p$.
\end{proof}

\begin{proposition}[Approximate Campbell identity]
\label{prop:Campbell}
For every family of functions
\[
 g_r:\Z_+^{\cI_n}\to[-1,1],
 \qquad r\in\cI_n,
\]
one has
\begin{align}
 &\left|
 \Ee\sum_{r\in\cI_n}C_{n,r}g_r(\xi_n-e_r)
 -\sum_{r\in\cI_n}a_{n,r}\Ee g_r(\xi_n)
 \right|
 \le\varepsilon_n,
 \label{eq:Campbell}
\end{align}
where
\begin{equation}\label{eq:epsilon-n}
 \varepsilon_n
 \le C_\kappa\varepsilon_n^{\mathrm{sw}}
 +C_\kappa(\log n)\left(\frac{U_n}{n}\right)^\beta
 \longrightarrow0.
\end{equation}
Here $g_r(\xi_n-e_r)$ is assigned any value when $C_{n,r}=0$, since it is
multiplied by $C_{n,r}$.
\end{proposition}

\begin{proof}
Fix $r\in\cI_n$.  Define a bounded function on permutations by
\[
 \Phi_r(\pi)=
 \begin{cases}
 g_r(\xi(\pi)-e_r),&C_r(\pi)\ge1,\\
 0,&C_r(\pi)=0,
 \end{cases}
\]
where $\xi(\pi)$ is the intermediate cycle vector of $\pi$.  Apply
\eqref{eq:marked-switch}.  Every switch on the right creates an $r$-cycle,
so the first case in the definition of $\Phi_r$ applies there.  On
$\{L_n(i)>r+U_n\}$, Lemma~\ref{lem:good-split} gives
\[
 \Phi_r\bigl(\sigma_n\circ(i\,\sigma_n^r(i))\bigr)
 =g_r(\xi_n).
\]
On the complementary source-cycle range, the difference between this
quantity and $g_r(\xi_n)$ is at most two.  Dividing the marked identity by
$r(n-r)$ therefore gives
\begin{align}
 &\left|
 \Ee[C_{n,r}g_r(\xi_n-e_r)]
 -a_{n,r}\Ee g_r(\xi_n)
 \right|
 \notag\\
 &\qquad\le
 a_{n,r}\Delta_{n,r}
 +2a_{n,r}
 \frac1n\Ee\sum_{i:L_n(i)>r}
 R_{i,\sigma_n^r(i)}
 \one_{\{L_n(i)\le r+U_n\}}.
 \label{eq:Campbell-one-r}
\end{align}

Since $r\le U_n\le n/2$, $a_{n,r}\le2/r$.  On a dyadic block
$a_j\le r<2a_j$,
\[
 \sum_{a_j\le r<2a_j}a_{n,r}\Delta_{n,r}
 \le C\sup_{a_j\le r<2a_j}\Delta_{n,r}.
\]
Proposition~\ref{prop:unweighted} handles the sum over blocks.  By
Lemma~\ref{lem:exceptional}, the second contribution summed over $r$ is at
most
\[
 C_\kappa\left(\frac{U_n}{n}\right)^\beta
 \sum_{r\in\cI_n}a_{n,r}
 \le C_\kappa(\log n)^{1-A_0\beta}=o(1),
\]
because $A_0\beta>1$.  This proves \eqref{eq:Campbell}.
\end{proof}

\section{A dimension-free product-Poisson Stein lemma}

We record the exact configuration-space statement needed to turn
Proposition~\ref{prop:Campbell} into total variation.

\begin{lemma}[Campbell-to-Poisson transfer]
\label{lem:Stein}
Let $I$ be finite, let $X=(X_i)_{i\in I}$ take values in $\Z_+^I$, and let
$\lambda_i\ge0$.  Suppose that for every family
$g_i:\Z_+^I\to[-1,1]$,
\begin{equation}\label{eq:abstract-Campbell}
 \left|
 \Ee\sum_{i\in I}X_i g_i(X-e_i)
 -\sum_{i\in I}\lambda_i\Ee g_i(X)
 \right|\le\varepsilon.
\end{equation}
Then
\begin{equation}\label{eq:Stein-TV}
 d_{\TV}\left(
 \cL(X),\bigotimes_{i\in I}\Pois(\lambda_i)
 \right)\le\varepsilon.
\end{equation}
The constant is independent of $|I|$ and of $\sum_i\lambda_i$.
\end{lemma}

\begin{proof}
Let $Z=(Z_i)_{i\in I}$ have the target product-Poisson law.  Consider the
immigration--death generator
\begin{align}
 \cA h(\eta)
 ={}&\sum_{i\in I}\lambda_i[h(\eta+e_i)-h(\eta)]
 \notag\\
 &+\sum_{i\in I}\eta_i[h(\eta-e_i)-h(\eta)].
 \label{eq:ID-generator}
\end{align}
For an indicator $f=\one_A$, let
\begin{equation}\label{eq:Stein-solution}
 h_f(\eta)=-\int_0^\infty
 \left(\Ee_\eta f(Y_t)-\Ee f(Z)\right)\dd t,
\end{equation}
where $(Y_t)$ is the immigration--death process with generator $\cA$.
Then
\[
 \cA h_f=f-\Ee f(Z).
\]
Couple processes started from $\eta$ and $\eta+e_i$ by using the same
immigrations and deaths.  Their only discrepancy is the extra initial
particle, which survives to time $t$ with probability $e^{-t}$.  Hence
\begin{equation}\label{eq:Stein-factor}
 \sup_{\eta,i}|h_f(\eta+e_i)-h_f(\eta)|
 \le\int_0^\infty e^{-t}\dd t=1.
\end{equation}
Put $g_i(\eta)=h_f(\eta+e_i)-h_f(\eta)$.  Then
\begin{align*}
 \Ee f(X)-\Ee f(Z)
 &=\Ee\cA h_f(X)\\
 &=\sum_i\lambda_i\Ee g_i(X)
 -\Ee\sum_iX_i g_i(X-e_i).
\end{align*}
Equation~\eqref{eq:abstract-Campbell} bounds the absolute value by
$\varepsilon$.  Take the supremum over events $A$.
\end{proof}

\begin{proof}[Proof of Theorem~\ref{thm:TV-main}]
Apply Lemma~\ref{lem:Stein} to Proposition~\ref{prop:Campbell} with
$\lambda_r=a_{n,r}$.  This gives
\begin{equation}\label{eq:TV-anr}
 d_{\TV}\left(
 \cL(\xi_n),
 \bigotimes_{r\in\cI_n}\Pois(a_{n,r})
 \right)\le\varepsilon_n=o(1).
\end{equation}
For Poisson variables,
\[
 d_{\TV}(\Pois(\lambda),\Pois(\mu))
 \le |\lambda-\mu|.
\]
Therefore
\begin{align*}
 d_{\TV}\left(
 \bigotimes_{r\in\cI_n}\Pois(a_{n,r}),
 \bigotimes_{r\in\cI_n}\Pois(1/r)
 \right)
 &\le\sum_{r\in\cI_n}\frac1{n-r}\\
 &\le C\frac{U_n}{n}=o(1).
\end{align*}
The triangle inequality proves \eqref{eq:TV-main}.
\end{proof}

\begin{remark}[Why the theorem is genuinely high-dimensional]
The number of coordinates is of order $U_n$, and the total Poisson intensity
is of order $\log n$.  The conclusion is not obtained by checking any fixed
collection of cycle lengths.  The dimension-free Stein factor in
\eqref{eq:Stein-factor} is what permits the summation of the one-switch
errors over the entire growing window.
\end{remark}

\section{The logarithmic Poisson field}

We first prove the functional theorem for the independent Poisson reference
process.  Let $(Z_{n,r})_{r\in\cI_n}$ be independent with
\begin{equation}\label{eq:Poisson-anr}
 Z_{n,r}\sim\Pois(a_{n,r}),
\end{equation}
and put
\[
 \ell_n=\log n,
 \qquad
 u_{n,r}=\frac{\log r}{\ell_n}.
\]
Define the two-dimensional process
\begin{equation}\label{eq:Mn}
 M_n(t)=\frac1{\sqrt{\ell_n}}
 \sum_{\substack{r\in\cI_n\\u_{n,r}\le t}}
 (Z_{n,r}-a_{n,r})
 \begin{pmatrix}1\\ \sqrt3\,u_{n,r}\end{pmatrix},
 \qquad 0\le t\le1.
\end{equation}

\subsection{Uniform logarithmic Riemann sums}

\begin{lemma}[Uniform logarithmic moment sums]
\label{lem:Riemann}
For $k=0,1,2$,
\begin{equation}\label{eq:Riemann}
 \sup_{0\le t\le1}
 \left|
 \frac1{\ell_n^{k+1}}
 \sum_{\substack{r\in\cI_n\\r\le n^t}}
 a_{n,r}(\log r)^k
 -\frac{t^{k+1}}{k+1}
 \right|
 \longrightarrow0.
\end{equation}
\end{lemma}

\begin{proof}
Use
\[
 a_{n,r}=\frac1r+\frac1{n-r}.
\]
The $1/(n-r)$ contribution is bounded uniformly in $t$ by
\[
 \frac{C}{\ell_n^{k+1}}
 \frac{U_n}{n}(\log n)^k=o(1).
\]
For the $1/r$ contribution, integral comparison gives, uniformly for
$2\le x\le n$,
\[
 \sum_{r\le x}\frac{(\log r)^k}{r}
 =\frac{(\log x)^{k+1}}{k+1}+O_k(1+(\log x)^k/x).
\]
Restricting to $b_n<r\le U_n$ removes at most
\[
 O((\log b_n)^{k+1})+O(\ell_n^k\log(n/U_n))
 =O(\ell_n^k\log\ell_n).
\]
After division by $\ell_n^{k+1}$ this is $o(1)$, uniformly in $t$.
\end{proof}

The deterministic covariance matrix of $M_n(t)$ is
\begin{equation}\label{eq:Qn}
 Q_n(t)=\frac1{\ell_n}
 \sum_{\substack{r\in\cI_n\\u_{n,r}\le t}}
 a_{n,r}
 \begin{pmatrix}
 1&\sqrt3u_{n,r}\\
 \sqrt3u_{n,r}&3u_{n,r}^2
 \end{pmatrix}.
\end{equation}
By Lemma~\ref{lem:Riemann},
\begin{equation}\label{eq:Q-limit}
 \sup_{0\le t\le1}\|Q_n(t)-Q(t)\|\longrightarrow0,
 \qquad
 Q(t)=
 \begin{pmatrix}
 t&\frac{\sqrt3}{2}t^2\\[1mm]
 \frac{\sqrt3}{2}t^2&t^3
 \end{pmatrix}.
\end{equation}

\subsection{A triangular-array functional CLT}

\begin{proposition}[Poisson functional CLT]
\label{prop:Poisson-FCLT}
The processes $M_n$ converge in $D([0,1],\R^2)$ to
\begin{equation}\label{eq:Gaussian-limit-process}
 G(t)=\left(W(t),\sqrt3\int_0^t u\,\dd W(u)\right).
\end{equation}
\end{proposition}

\begin{proof}
Order the independent centered increments in \eqref{eq:Mn} by increasing
$r$.  Their covariance partial sums converge uniformly to $Q(t)$ by
\eqref{eq:Q-limit}.  It remains to verify the Lindeberg condition.

Let
\[
 v_{n,r}=\frac1{\sqrt{\ell_n}}
 \begin{pmatrix}1\\ \sqrt3u_{n,r}\end{pmatrix}.
\]
Then $\|v_{n,r}\|\le2/\sqrt{\ell_n}$.  If
$Z\sim\Pois(\lambda)$, then
\[
 \Ee(Z-\lambda)^4=\lambda+3\lambda^2.
\]
Consequently,
\begin{align*}
 \sum_{r\in\cI_n}
 \Ee\left\|v_{n,r}(Z_{n,r}-a_{n,r})\right\|^4
 &\le\frac{C}{\ell_n^2}
 \sum_{r\in\cI_n}(a_{n,r}+a_{n,r}^2)\\
 &=O(\ell_n^{-1}).
\end{align*}
For every $\varepsilon>0$, the Lindeberg sum is bounded by
$\varepsilon^{-2}$ times this fourth-moment sum and therefore tends to zero.
The standard martingale or independent-triangular-array functional CLT now
gives a continuous centered Gaussian martingale with covariance $Q$.
Since
\[
 Q(t)=\int_0^t
 \begin{pmatrix}1\\ \sqrt3u\end{pmatrix}
 \begin{pmatrix}1&\sqrt3u\end{pmatrix}\dd u,
\]
the limit has the representation \eqref{eq:Gaussian-limit-process}.
\end{proof}

\begin{remark}[Continuous-limit tightness]
The fourth-moment calculation also implies that the maximal jump tends to
zero in probability.  Hence the sequence is $C$-tight, although we state the
result in the standard $J_1$ topology on $D([0,1],\R^2)$.
\end{remark}

\section{Transfer to the Luce intermediate field}

Define
\begin{equation}\label{eq:intermediate-process}
 \widehat M_n(t)=\frac1{\sqrt{\ell_n}}
 \sum_{\substack{r\in\cI_n\\u_{n,r}\le t}}
 (C_{n,r}-a_{n,r})
 \begin{pmatrix}1\\ \sqrt3u_{n,r}\end{pmatrix}.
\end{equation}

\begin{proposition}[Intermediate Luce functional CLT]
\label{prop:Luce-intermediate-FCLT}
One has
\begin{equation}\label{eq:Luce-intermediate-FCLT}
 \widehat M_n\Rightarrow
 \left(W(t),\sqrt3\int_0^t u\,\dd W(u)\right)
\end{equation}
in $D([0,1],\R^2)$.
\end{proposition}

\begin{proof}
The map from the finite vector $(x_r)_{r\in\cI_n}$ to the path in
\eqref{eq:intermediate-process} is measurable.  Total variation cannot
increase under a measurable map.  Theorem~\ref{thm:TV-main} therefore gives
\[ d_{\TV}(\cL(\widehat M_n),\cL(M_n))\longrightarrow0.
\]
Apply Proposition~\ref{prop:Poisson-FCLT}.
\end{proof}

The same argument gives arbitrary finite families of bounded logarithmic
weights.

\begin{proposition}[Intermediate Gaussian field]
\label{prop:intermediate-field}
Let $f_1,\ldots,f_d$ be bounded Riemann-integrable functions on $[0,1]$.
Then
\begin{align}
 &\left(
 \frac1{\sqrt{\ell_n}}
 \sum_{r\in\cI_n}f_j(u_{n,r})(C_{n,r}-a_{n,r})
 \right)_{j=1}^d
 \notag\\
 &\hspace{2cm}\Rightarrow
 N(0,\Sigma),
 \qquad
 \Sigma_{ij}=\int_0^1f_i(u)f_j(u)\,\dd u.
 \label{eq:intermediate-field}
\end{align}
\end{proposition}

\begin{proof}
Transfer to the product-Poisson field by total variation.  For any fixed
linear combination $f=\sum_jt_jf_j$, the centered Poisson sum has variance
\[
 \frac1{\ell_n}\sum_{r\in\cI_n}a_{n,r}f(u_{n,r})^2
 \longrightarrow\int_0^1f(u)^2\dd u.
\]
The maximal coefficient is $O(\ell_n^{-1/2})$, and the same fourth-moment
argument as above gives Lindeberg.  Apply Cram\'er--Wold.
\end{proof}

\section{Restoring the full cycle processes}

Define the intermediate count and additive processes
\begin{align}
 K_n^\circ(t)&=\sum_{\substack{r\in\cI_n\\r\le n^t}}C_{n,r},
 \label{eq:Kcirc}\\
 B_n^\circ(t)&=\sum_{\substack{r\in\cI_n\\r\le n^t}}(\log r)C_{n,r}.
 \label{eq:Bcirc-process}
\end{align}
Because all omitted summands are nonnegative,
\begin{align}
 \sup_{0\le t\le1}|K_n(t)-K_n^\circ(t)|
 &\le\sum_{r\le b_n}C_{n,r}+\sum_{r>U_n}C_{n,r},
 \label{eq:K-uniform-omit}\\
 \sup_{0\le t\le1}|B_n(t)-B_n^\circ(t)|
 &\le\sum_{r\le b_n}(\log r)C_{n,r}
 +\sum_{r>U_n}(\log r)C_{n,r}.
 \label{eq:B-uniform-omit}
\end{align}

\begin{lemma}[Uniform restoration of omitted cycle lengths]
\label{lem:restore}
One has
\begin{align}
 \sup_{0\le t\le1}
 \frac{|K_n(t)-K_n^\circ(t)|}{\sqrt{\ell_n}}
 &\longrightarrow0
 &&\text{in }L^1,
 \label{eq:restore-K}\\
 \sup_{0\le t\le1}
 \frac{|B_n(t)-B_n^\circ(t)|}{\ell_n^{3/2}}
 &\longrightarrow0
 &&\text{in }L^1.
 \label{eq:restore-B}
\end{align}
\end{lemma}

\begin{proof}
By Lemma~\ref{lem:crude}, the expectation of the right side in
\eqref{eq:K-uniform-omit} is
\[
 O_\kappa(\log b_n+1+\log(n/U_n))
 =O_\kappa(\log\ell_n)=o(\sqrt{\ell_n}).
\]
The expectation of the right side in \eqref{eq:B-uniform-omit} is
\[
 O_\kappa((\log b_n)^2+\ell_n\{1+\log(n/U_n)\})
 =O_\kappa(\ell_n\log\ell_n)
 =o(\ell_n^{3/2}).
\]
\end{proof}

We next replace the exact intermediate Poisson centers by the classical
centers.

\begin{lemma}[Uniform deterministic centering]
\label{lem:centering}
Uniformly for $0\le t\le1$,
\begin{align}
 \sum_{\substack{r\in\cI_n\\r\le n^t}}a_{n,r}
 &=t\ell_n+O(\log\ell_n),
 \label{eq:center-K}\\
 \sum_{\substack{r\in\cI_n\\r\le n^t}}a_{n,r}\log r
 &=\frac12t^2\ell_n^2+O(\ell_n\log\ell_n).
 \label{eq:center-B}
\end{align}
\end{lemma}

\begin{proof}
This is Lemma~\ref{lem:Riemann} for $k=0,1$, with the error from deleting the
lower and upper ranges retained rather than divided by the leading scale.
The $1/(n-r)$ portion contributes $O(U_n/n)$ in \eqref{eq:center-K} and
$O(\ell_nU_n/n)$ in \eqref{eq:center-B}.
\end{proof}

\begin{theorem}[Bivariate additive functional theorem]
\label{thm:additive-functional}
In $D([0,1],\R^2)$,
\begin{equation}\label{eq:additive-functional}
 \boxed{
 \left(
 \frac{K_n(t)-t\ell_n}{\sqrt{\ell_n}},
 \frac{B_n(t)-\frac12t^2\ell_n^2}{\sqrt{\ell_n^3/3}}
 \right)_{0\le t\le1}
 \Rightarrow
 \left(
 W(t),\sqrt3\int_0^t u\,\dd W(u)
 \right)_{0\le t\le1}.}
\end{equation}
\end{theorem}

\begin{proof}
The two coordinates of \eqref{eq:intermediate-process} are
\[
 \frac{K_n^\circ(t)-\sum_{r\in\cI_n,r\le n^t}a_{n,r}}
      {\sqrt{\ell_n}}
\]
and
\[
 \frac{B_n^\circ(t)-\sum_{r\in\cI_n,r\le n^t}a_{n,r}\log r}
      {\sqrt{\ell_n^3/3}}.
\]
Apply Proposition~\ref{prop:Luce-intermediate-FCLT}, then use
Lemmas~\ref{lem:restore} and \ref{lem:centering} and Slutsky's theorem in
$D([0,1],\R^2)$.
\end{proof}

\begin{proof}[Proof of Theorem~\ref{thm:isonormal}]
Apply Proposition~\ref{prop:intermediate-field}.  Replacing $a_{n,r}$ by
$1/r$ changes each normalized deterministic center by
\[
 O\left(\frac1{\sqrt{\ell_n}}\sum_{r\le U_n}\frac1{n-r}\right)=o(1).
\]
For a bounded $f$, the omitted random contribution is bounded by
$\|f\|_\infty$ times the omitted cycle count in
\eqref{eq:K-uniform-omit}, which is $o_{L^1}(\sqrt{\ell_n})$.  The omitted
$1/r$ center has size $O(\log\ell_n)$.  Cram\'er--Wold completes the proof.
\end{proof}

\section{Functional additive-to-LCM transfer}

For $m\ge1$, define the cumulative divisibility count
\begin{equation}\label{eq:Dmt}
 D_{n,m}(t)=\sum_{\substack{r\le n^t\\m\mid r}}C_{n,r}.
\end{equation}
Let $\Lambda$ be the von Mangoldt function.  The identities
\[
 \log r=\sum_{m\mid r}\Lambda(m)
\]
and
\[
 \cO_n(t)=\sum_{m\le n^t}\Lambda(m)
 \one_{\{D_{n,m}(t)\ge1\}}
\]
give the exact process decomposition
\begin{equation}\label{eq:process-vonmangoldt}
 B_n(t)-\cO_n(t)
 =\sum_{m\le n^t}\Lambda(m)(D_{n,m}(t)-1)_+.
\end{equation}

\begin{lemma}[Monotonicity of the product--LCM discrepancy]
\label{lem:monotone-discrepancy}
The process
\[
 t\longmapsto B_n(t)-\cO_n(t)
\]
is nonnegative and nondecreasing.  Consequently,
\begin{equation}\label{eq:sup-discrepancy}
 \sup_{0\le t\le1}|B_n(t)-\cO_n(t)|
 =B_n(1)-\log\ord(\sigma_n).
\end{equation}
\end{lemma}

\begin{proof}
For every $m$, $D_{n,m}(t)$ is nondecreasing in $t$, and
$x\mapsto(x-1)_+$ is nondecreasing.  Apply
\eqref{eq:process-vonmangoldt}.
\end{proof}

\begin{proposition}[Uniform functional LCM transfer]
\label{prop:functional-LCM-transfer}
One has
\begin{equation}\label{eq:functional-LCM-transfer}
 \sup_{0\le t\le1}
 \frac{|B_n(t)-\cO_n(t)|}{\ell_n^{3/2}}
 \longrightarrow0
\end{equation}
in $L^1$.
\end{proposition}

\begin{proof}
By Lemma~\ref{lem:monotone-discrepancy} and
Input Theorem~\ref{input:arithmetic},
\[
 \Ee\sup_{0\le t\le1}|B_n(t)-\cO_n(t)|
 =\Ee[B_n-\log O_n]
 \le C_\kappa\ell_n\log\ell_n
 =o(\ell_n^{3/2}).
\]
\end{proof}

\begin{proof}[Proof of Theorem~\ref{thm:main-functional}]
After a common rescaling, the rate hypothesis is exactly
\eqref{eq:comparable}.  Theorem~\ref{thm:additive-functional} gives the
bivariate count/additive-process limit.  Proposition
\ref{prop:functional-LCM-transfer} replaces $B_n(t)$ by $\cO_n(t)$ uniformly
in $t$ at the second-coordinate scale.  Slutsky's theorem proves
\eqref{eq:main-functional}.
\end{proof}

\section{Multivariate consequences}

The process theorem immediately gives all fixed finite-time joint laws.

\begin{corollary}[Finite-time multivariate law]
\label{cor:finite-time}
Fix $0\le t_1<\cdots<t_d\le1$.  The $2d$-vector consisting of
\[
 \frac{K_n(t_j)-t_j\ell_n}{\sqrt{\ell_n}}
 \quad\text{and}\quad
 \frac{\cO_n(t_j)-\frac12t_j^2\ell_n^2}
      {\sqrt{\ell_n^3/3}},
 \qquad 1\le j\le d,
\]
converges to a centered Gaussian vector whose covariance is determined by
\eqref{eq:cov-count-count}--\eqref{eq:cov-order-order}.
\end{corollary}

\begin{proof}
Apply the continuous evaluation map at the fixed times to
Theorem~\ref{thm:main-functional}; the limit paths are continuous.
\end{proof}

\begin{proof}[Proof of Corollary~\ref{cor:endpoint-joint}]
Take $t=1$ in Theorem~\ref{thm:main-functional}.  Equations
\eqref{eq:cov-count-count}--\eqref{eq:cov-order-order} give the displayed
covariance matrix.
\end{proof}

\begin{corollary}[Independent logarithmic annuli]
\label{cor:annuli}
Fix disjoint intervals $(s_j,t_j]\subset[0,1]$.  The corresponding normalized
increment vectors
\[
 \left(
 \frac{K_n(t_j)-K_n(s_j)-(t_j-s_j)\ell_n}{\sqrt{\ell_n}},
 \frac{\cO_n(t_j)-\cO_n(s_j)-\frac12(t_j^2-s_j^2)\ell_n^2}
      {\sqrt{\ell_n^3/3}}
 \right)
\]
converge jointly to independent centered Gaussian vectors.  For one annulus
$(s,t]$, the covariance matrix is
\[
 \begin{pmatrix}
 t-s&\frac{\sqrt3}{2}(t^2-s^2)\\[1mm]
 \frac{\sqrt3}{2}(t^2-s^2)&t^3-s^3
 \end{pmatrix}.
\]
\end{corollary}

\begin{proof}
The limiting increments are stochastic integrals against $W$ over disjoint
intervals and are therefore independent.
\end{proof}

\begin{corollary}[Time-changed Brownian order process]
\label{cor:time-change}
The second coordinate of Theorem~\ref{thm:main-functional} may be written as
\[
 \left(B(t^3)\right)_{0\le t\le1}
\]
in distribution, where $B$ is standard Brownian motion.
\end{corollary}

\begin{proof}
The process $\sqrt3\int_0^tu\,\dd W(u)$ is a centered continuous Gaussian
martingale with quadratic variation
\[
 3\int_0^t u^2\dd u=t^3.
\]
Apply the Dambis--Dubins--Schwarz representation, or compare Gaussian
covariances directly.
\end{proof}

At the endpoint, the common-noise representation gives a useful regression
identity.

\begin{corollary}[Endpoint Gaussian regression]
\label{cor:regression}
Let $(Z_K,Z_O)$ denote the limiting endpoint pair.  Then
\[
 Z_O\stackrel d=\frac{\sqrt3}{2}Z_K+\frac12Z_\perp,
\]
where $Z_K,Z_\perp$ are independent standard normals.  Thus the normalized
cycle count explains three quarters of the limiting order variance.
\end{corollary}

\section{Relation to the Berry--Esseen branch}

The product-Poisson theorem is a qualitative process tool.  It should not be
confused with the frequency-sensitive comparison used in the quantitative
supplement \cite{LongBerry2026}.

The total-variation error in \eqref{eq:epsilon-n} contains the reset term
\[
 \exp\{-c_\kappa\sqrt{\log\log n}\},
\]
up to polylogarithmic factors.  This tends to zero, but it is much larger
than the natural scalar Berry--Esseen scale $(\log n)^{-1/2}$.  In the scalar
characteristic-function ODE, the same raw switching error is divided by the
standard deviation and integrated over a controlled frequency window.  The
Berry--Esseen supplement exploits that extra structure to obtain the sharp
$O_\kappa((\log n)^{-1/2})$ rate for the additive proxy.

Accordingly, the recommended division of labor is:

\begin{enumerate}[label=\textnormal{\arabic*.},leftmargin=2.2em]
\item Use Theorem~\ref{thm:TV-main} and the present branch for qualitative
functional and multivariate universality.

\item Retain the scalar switching ODE for sharp additive Berry--Esseen and
Edgeworth analysis.

\item Use the multi-clock moment bounds from the quantitative supplement for
concentration of the arithmetic remainder and order-level rates.

\item Pursue a quantitative functional theorem only after selecting an
explicit path metric.  A sharp endpoint Kolmogorov rate and a path-space
bounded-Lipschitz rate are different problems and need not have the same
proof.
\end{enumerate}

\begin{remark}[Why this branch still helps quantitative work]
The product-Poisson approximation identifies the correct Gaussian covariance
operator and proves that no hidden profile-dependent covariance survives on
logarithmic scale.  It therefore supplies the structural target for any
future multivariate Stein, strong approximation, or functional
Berry--Esseen theorem, even though its present total-variation error is not
sharp enough for the scalar rate.
\end{remark}

\section{Audit checklist and merge plan}

Before this branch replaces the scalar section in the main paper, the
following points should be checked independently.

\begin{enumerate}[label=\textnormal{(\roman*)},leftmargin=2.4em]
\item Re-run Proposition~\ref{prop:unweighted} directly from the exact
constants in the finite-endpoint input, with the stronger choice
$A_0\beta>1$ recorded explicitly.

\item Verify the permutation convention in Lemma~\ref{lem:good-split} against
the marked switch used in the companion manuscript: right-composition by
$(i\,\sigma_n^r(i))$ must split the source cycle into lengths $r$ and
$L-r$ with no orientation reversal.

\item Check the factor two and normalization $a_{n,r}=n/[r(n-r)]$ in
\eqref{eq:Campbell-one-r}.

\item Keep the Stein lemma in the paper, rather than relying on an unstated
high-dimensional Poisson approximation theorem.  Its dimension-free first
difference bound is the central reason the argument closes.

\item Verify the uniform-in-$t$ deterministic center estimates at the two
cutoff transition times $t=\log b_n/\log n$ and
$t=\log U_n/\log n$.

\item Preserve the scalar characteristic-function proof in a separate branch
until the product-Poisson theorem has passed this audit.  Afterward, either
move the scalar proof to an appendix or retain it only in the Berry--Esseen
supplement.
\end{enumerate}

A merged final paper would naturally have the following structure:

\begin{enumerate}[label=\textnormal{\arabic*.},leftmargin=2.2em]
\item Introduction and the bivariate functional theorem.
\item Model, marked one-switch input, and quantitative endpoint interface.
\item Unweighted dyadic summability.
\item Product-Poisson approximation of the intermediate cycle field.
\item Gaussian random measure and functional additive laws.
\item Rate-oriented three-clock switch and factorial moments.
\item Functional product--LCM transfer.
\item Scalar, joint, and annular consequences.
\item Quantitative directions and relation to the Berry--Esseen supplement.
\end{enumerate}

\appendix

\section{A general weighted additive-to-LCM process transfer}

The monotonicity argument is independent of permutations and may be useful in
other component models.

Let $(C_{n,r})_{r\ge1}$ be arbitrary nonnegative integer-valued variables and
let $x_n(t)$ be a nondecreasing integer-valued cutoff.  Define
\[
 A_n(t)=\sum_{r\le x_n(t)}C_{n,r}\log r,
\]
and
\[
 L_n(t)=\log\lcm\{r\le x_n(t):C_{n,r}>0\}.
\]

\begin{proposition}[Abstract functional additive-to-LCM transfer]
\label{prop:abstract-transfer}
The discrepancy $A_n(t)-L_n(t)$ is nonnegative and nondecreasing.  Hence, for
any deterministic scale $s_n>0$,
\[
 \sup_t\frac{|A_n(t)-L_n(t)|}{s_n}
 =\frac{A_n(1)-L_n(1)}{s_n}.
\]
Consequently, if
\[
 \frac{A_n(\cdot)-a_n(\cdot)}{s_n}\Rightarrow X(\cdot)
\]
in $D$ and
\[
 \frac{A_n(1)-L_n(1)}{s_n}\longrightarrow0
\]
in probability, then
\[
 \frac{L_n(\cdot)-a_n(\cdot)}{s_n}\Rightarrow X(\cdot).
\]
\end{proposition}

\begin{proof}
Put
\[
 D_{n,m}(t)=\sum_{\substack{r\le x_n(t)\\m\mid r}}C_{n,r}.
\]
Then
\[
 A_n(t)-L_n(t)
 =\sum_m\Lambda(m)(D_{n,m}(t)-1)_+.
\]
Each summand is nonnegative and nondecreasing in $t$.  The functional transfer
is Slutsky's theorem under the uniform perturbation displayed above.
\end{proof}

\section{A finite-dimensional covariance formula}

For later quantitative work, it is useful to record the covariance of all
linear combinations of the two limiting processes.  Let
\[
 X(t)=W(t),
 \qquad
 Y(t)=\sqrt3\int_0^t u\,\dd W(u).
\]
For coefficients $a_j,b_j$ and times $t_j$,
\[
 \sum_j\{a_jX(t_j)+b_jY(t_j)\}
 =\int_0^1
 \left[
 \sum_j\{a_j+\sqrt3b_ju\}\one_{\{u\le t_j\}}
 \right]\dd W(u).
\]
Therefore its variance is
\begin{equation}\label{eq:general-variance}
 \int_0^1
 \left[
 \sum_j\{a_j+\sqrt3b_ju\}\one_{\{u\le t_j\}}
 \right]^2\dd u.
\end{equation}
This is the covariance form that should be preserved in any multivariate
Berry--Esseen or Edgeworth extension.

\begin{thebibliography}{99}

\bibitem{ABT2003}
R. Arratia, A. D. Barbour, and S. Tavar\'e,
\newblock \emph{Logarithmic Combinatorial Structures: A Probabilistic
Approach},
\newblock EMS Monographs in Mathematics, European Mathematical Society,
Z\"urich, 2003.

\bibitem{BarbourHolstJanson1992}
A. D. Barbour, L. Holst, and S. Janson,
\newblock \emph{Poisson Approximation},
\newblock Oxford Studies in Probability, vol.~2, Clarendon Press, Oxford,
1992.

\bibitem{BarbourTavare1994}A. D. Barbour and S. Tavar\'e,
\newblock A rate for the Erd\H{o}s--Tur\'an law,
\newblock \emph{Combinatorics, Probability and Computing} 3 (1994), no.~2,
167--176.

\bibitem{Billingsley1999}
P. Billingsley,
\newblock \emph{Convergence of Probability Measures}, second edition,
\newblock Wiley Series in Probability and Statistics, Wiley, New York, 1999.

\bibitem{CDK2024}
S. Chatterjee, P. Diaconis, and G. B. Kim,
\newblock Enumerative theory for the Tsetlin library,
\newblock \emph{Journal of Algebra} 655 (2024), 139--162.

\bibitem{ErdosTuran1967}
P. Erd\H{o}s and P. Tur\'an,
\newblock On some problems of a statistical group theory. III,
\newblock \emph{Acta Mathematica Academiae Scientiarum Hungaricae} 18
(1967), 309--320.

\bibitem{EthierKurtz1986}
S. N. Ethier and T. G. Kurtz,
\newblock \emph{Markov Processes: Characterization and Convergence},
\newblock Wiley Series in Probability and Mathematical Statistics, Wiley,
New York, 1986.

\bibitem{LongBerry2026}
C. D. Long,
\newblock Berry--Esseen checkpoints for the Erd\H{o}s--Tur\'an law of
comparable-rate Luce permutations,
\newblock research supplement, 2026.

\bibitem{LongOrder2026}
C. D. Long,
\newblock An Erd\H{o}s--Tur\'an law for comparable-rate Luce permutations,
\newblock manuscript, 2026.

\bibitem{LongSwitch2026}
C. D. Long,
\newblock Mesoscopic cycles of comparable-rate Luce permutations: a
switching proof under submacroscopic root-cycle escape,
\newblock manuscript, 2026.

\bibitem{Luce1959}
R. D. Luce,
\newblock \emph{Individual Choice Behavior: A Theoretical Analysis},
\newblock John Wiley \& Sons, New York, 1959.

\end{thebibliography}

\end{document}